\chapter{Change des marchés émergents et NDF}\label{ch:m2:emerging-market-fx-and-ndfs}

Pendant trois ans et quatre mois, l'euro ne valut jamais moins de 1.20 franc
suisse. Le 6 septembre 2011, la Banque nationale suisse avait déclaré
qu'elle «~ne tolérerait plus~» un taux inférieur et qu'elle était «~prête à acheter des devises en
quantités illimitées~», et le taux de référence quotidien de la Banque centrale européenne,
1.1111 la veille, ne retomba plus jamais sous 1.20. Le 15 janvier
2015, la Banque nationale suisse abandonna ce plancher. À l'heure de référence de la Banque centrale
européenne, cet après-midi-là, l'euro valait 1.0280 franc, soit 14.4~\%
de moins que la veille, et une semaine plus tard 0.9816. Un courtier en devises en ligne
new-yorkais constata que ses clients lui devaient davantage que leurs dépôts et emprunta
300~millions de dollars en un jour pour rester en activité. Les taux de change
qu'un gouvernement fixe, encadre dans une bande ou contrôle, restent calmes pendant des années, puis cessent de l'être.
Ce chapitre traite des monnaies des marchés émergents, et de quelques monnaies de pays riches,
où l'État se tient entre l'acheteur et le vendeur~: convertibilité
et \omterm{def:m2:emerging-market-fx-and-ndfs:controls}{contrôles des capitaux}, les \omterm{def:m2:emerging-market-fx-and-ndfs:onshore}{marchés onshore} et offshore qu'ils font naître, le
\omterm{def:m2:emerging-market-fx-and-ndfs:ndf}{contrat à terme non livrable} que le \omterm{def:m2:emerging-market-fx-and-ndfs:onshore}{marché offshore} négocie, ainsi que les ancrages, les bandes et le
jour de leur rupture.

\section{Convertibilité et contrôles des capitaux}

\begin{definition}[Convertibilité, contrôle des capitaux]\label{def:m2:emerging-market-fx-and-ndfs:controls}
Une devise est \emph{convertible}\index{convertible} lorsque toute personne peut l'échanger
contre une devise étrangère, quel qu'en soit l'usage, au taux du marché. Un \emph{contrôle des
capitaux}\index{contrôle des capitaux} est une restriction légale portant sur de tels échanges ou sur
la circulation transfrontalière de l'argent~: limites sur qui peut acheter des devises étrangères, dans
quel but, pour quel montant, ou à quelle vitesse elles peuvent sortir du pays ou y entrer.
\end{definition}

La plupart des devises des pays riches sont entièrement \omterm{def:m2:emerging-market-fx-and-ndfs:controls}{convertibles}. Beaucoup de devises émergentes
sont \omterm{def:m2:emerging-market-fx-and-ndfs:controls}{convertibles} pour les échanges commerciaux~: une entreprise qui paie des importations peut acheter
des dollars, mais pas librement pour investir~: un investisseur étranger peut devoir déclarer
ses achats d'obligations locales, et un résident local peut être limité quant au montant
qu'il peut détenir à l'étranger. Ces contrôles existent pour que le taux de change et le
crédit intérieur restent sous le contrôle des autorités, et ils ont une conséquence directe pour
les marchés~: la devise ne peut pas être livrée hors du pays, de sorte que les étrangers qui
veulent se couvrir ou spéculer ont besoin d'un instrument qui se règle dans une autre monnaie.

\section{Marchés onshore et offshore}

\begin{definition}[Marché onshore, marché offshore]\label{def:m2:emerging-market-fx-and-ndfs:onshore}
Le \emph{marché onshore}\index{marché onshore} d'une devise est le marché de
son propre pays, sous l'autorité de ses régulateurs et de ses contrôles~; le \emph{marché
offshore}\index{marché offshore} est le marché situé hors du pays, entre non-résidents,
où la devise se négocie à d'autres conditions ou seulement sous forme réglée en espèces.
\end{definition}

Le renminbi illustre les deux à la fois. En onshore, le renminbi (code CNY) se négocie sur le
marché interbancaire autour d'un cours central que le China Foreign Exchange Trade
System publie chaque matin pour la Banque populaire de Chine~; depuis le 17 mars
2014, le cours au comptant face au dollar peut s'écarter de 2~\% au-dessus ou au-dessous de ce cours. Hong
Kong, premier centre offshore, négocie la même devise en dehors de ces règles,
sous le code CNH, avec ses propres prix au comptant et à terme et son propre fixing, publié par
la Treasury Markets Association. Les deux cours restent en général proches, car
les flux entre les marchés sont autorisés par des canaux contrôlés, et ils divergent
lorsque les contrôles deviennent contraignants (voir la \cref{fig:m2:emerging-market-fx-and-ndfs:markets}).
Leurs contrats à terme impliquent des taux d'intérêt différents pour le renminbi~: une
base onshore--offshore.

\begin{omfigure}
\centering
\begin{tikzpicture}[font=\small,
  box/.style={draw,thick,rounded corners=2pt,align=center,minimum height=1.2cm,text width=4.6cm,inner sep=4pt},
  arr/.style={{Stealth[length=2.5mm]}-{Stealth[length=2.5mm]},thick}]
\node[box,draw=omThm,fill=omThm!8] (on) at (0,0) {\textbf{onshore (CNY)}\\marché interbancaire, résidents~;\\comptant à moins de 2~\% du\\cours central quotidien};
\node[box,draw=omDef,fill=omDef!8] (off) at (7.2,0) {\textbf{offshore (CNH)}\\Hong Kong et autres centres~;\\renminbi livrable,\\comptant, terme et fixing propres};
\node[box,draw=omProp,fill=omProp!8] (ndf) at (7.2,-2.6) {\textbf{NDF}\\réglés en espèces en dollars\\sur la base du fixing onshore};
\draw[arr,dashed] (on) -- node[above,font=\footnotesize,align=center] {canaux\\contrôlés} (off);
\draw[omProp,thick,-{Stealth[length=2.5mm]}] (on.south) |- node[pos=0.75,below,font=\footnotesize] {fixing du CNY (SAEC)} (ndf.west);
\end{tikzpicture}

\omcaption{Une devise, trois marchés. Le renminbi onshore se négocie dans une bande
autour du cours central quotidien~; le renminbi offshore est livrable hors de
Chine continentale à ses propres prix~; les NDF se règlent en dollars sur la base du
fixing onshore. Les fonds ne circulent entre les deux premiers qu'à travers des canaux contrôlés, de
sorte que leurs prix peuvent différer. Schéma.}\label{fig:m2:emerging-market-fx-and-ndfs:markets}
\end{omfigure}

\begin{example}[Deux courbes à terme pour une même devise]\label{ex:m2:emerging-market-fx-and-ndfs:basis}
Supposons que l'USDCNY onshore se négocie à 7.1000 au comptant et à 7.0400 à trois mois,
et l'USDCNH offshore à 7.1200 et 7.0700, avec un taux en dollars de 3.68~\%. Le
terme onshore implique un taux pour le renminbi de 0.31~\% sur 91 jours, celui de l'offshore
0.89~\%~: une base de 0.58 point de pourcentage, soit le prix du passage du renminbi de part et d'autre de la
frontière. La bande onshore ce jour-là irait de 6.9580 à 7.2420.
\end{example}

\section{Contrats à terme non livrables}

\begin{definition}[Contrat à terme non livrable, source du fixing]\label{def:m2:emerging-market-fx-and-ndfs:ndf}
Un \emph{contrat à terme non livrable}\index{contrat à terme non livrable} (NDF) est un
contrat à terme sur une paire de devises qui se règle en espèces, en dollars, à hauteur de l'écart entre le
taux contractuel $K$ et un taux de fixing publié $S$ à une date de
fixing~: pour un notionnel en dollars $N$, l'acheteur de dollars reçoit
$N(S - K)/S$. La \emph{source du fixing}\index{source du fixing} est le taux que le
contrat désigne pour $S$, généralement un taux publié par la banque centrale locale ou
par un agent officiel.
\end{definition}

Le NDF permet à un investisseur étranger de se couvrir, ou à un fonds de spéculer, sur une devise
qu'il ne peut ni détenir ni livrer~: aucune monnaie locale ne change de mains, seulement des dollars, hors
du pays. Son prix suit le fixing attendu et le taux du dollar, et
implique donc un taux d'intérêt offshore qui peut différer du taux onshore de la
base de l'\cref{ex:m2:emerging-market-fx-and-ndfs:basis}. Les
conditions standard du secteur désignent la \omterm{def:m2:emerging-market-fx-and-ndfs:ndf}{source du fixing} pour chaque devise et fixent le cours
deux jours ouvrés avant le règlement-livraison, comme le montre le tableau ci-dessous.

\begin{center}
\small
\begin{tabular}{@{}lll@{}}
\toprule
Devise & \omterm{def:m2:emerging-market-fx-and-ndfs:ndf}{Source du fixing} (code EMTA) & Publiée par, vers \\
\midrule
Roupie indienne & INR FBIL & Financial Benchmarks India, 13:30 Mumbai \\
Real brésilien & BRL PTAX & Banco Central do Brasil, 13:15 S\~ao Paulo \\
Renminbi chinois & CNY SAEC & CFETS pour la Banque populaire, 9:15 Pékin \\
Won coréen & KRW KFTC18 & Seoul Money Brokerage, 16:00 Séoul \\
Dollar de Taïwan & TWD TAIFX1 & Taipei Forex, 11:00 Taipei \\
\bottomrule
\end{tabular}
\par\smallskip
{\footnotesize Sources de fixing de cinq devises faisant l'objet de NDF, chacune «~pour un règlement-livraison à deux
jours ouvrés~», dans l'annexe A des 1998 FX and Currency Option Definitions (à jour en
juin 2023).}
\end{center}

\begin{omfigure}
\centering
\begin{tikzpicture}[font=\small,x=1.35cm]
\draw[thick,-{Stealth}] (0,0) -- (9.6,0) node[right,font=\footnotesize] {temps};
\foreach \x/\l in {0/transaction~: $K$ convenu,6.4/{date de fixing~:\\$S$ publié},8.4/règlement} {\draw (\x,0.12) -- (\x,-0.12); \node[below,font=\footnotesize,align=center] at (\x,-0.15) {\l};}
\draw[gray,decorate,decoration={brace,amplitude=4pt}] (6.4,0.35) -- (8.4,0.35) node[midway,above=4pt,font=\footnotesize,gray] {deux jours ouvrés};
\draw[omThm,thick,-{Stealth}] (8.4,-1.2) -- (8.4,-2.0);
\node[omThm,font=\footnotesize,align=center] at (8.4,-2.45) {un paiement en dollars~:\\$N(S-K)/S$ à l'acheteur de dollars};
\node[omDef,font=\footnotesize,align=center] at (3.2,-1.6) {aucune devise locale ne change de mains~;\\rien n'a besoin d'entrer ni de sortir du pays};
\end{tikzpicture}

\omcaption{Le cycle de vie d'un \omterm{def:m2:emerging-market-fx-and-ndfs:ndf}{contrat à terme non livrable} (NDF). Le taux du contrat est fixé à la
conclusion~; deux jours ouvrés avant le règlement-livraison, la \omterm{def:m2:emerging-market-fx-and-ndfs:ndf}{source du fixing} publie le
taux, et à la date de règlement la différence est versée en dollars.
Schéma.}\label{fig:m2:emerging-market-fx-and-ndfs:ndf}
\end{omfigure}

\begin{example}[Un NDF sur le won]\label{ex:m2:emerging-market-fx-and-ndfs:krw}
Un fonds achète à terme trois mois USD~10 millions contre won dans un NDF à 1\,380.00,
avec règlement-livraison le mardi 13 octobre 2026. Le fixing est relevé deux jours ouvrés
avant, le vendredi 9 octobre, à 1\,425.50~: le fonds reçoit
$10\,000\,000 \times 45.50/1\,425.50 = \text{USD}~319\,186$. Si le won avait été fixé à
1\,350.00, le fonds aurait versé USD~222\,222.
\end{example}

\section{Ancrages, bandes et interventions}

\begin{definition}[Ancrage de change, bande de fluctuation, intervention de change]\label{def:m2:emerging-market-fx-and-ndfs:peg}
Un \emph{ancrage de change}\index{ancrage de change} est une politique consistant à maintenir le taux
de change à un niveau fixe par rapport à une autre devise~; une \emph{bande de
fluctuation}\index{bande de fluctuation} laisse ce taux évoluer dans des limites fixées autour
d'un taux central. Une \emph{intervention de change}\index{intervention de change} est l'achat ou la vente
de devises étrangères par les autorités pour faire évoluer ou maintenir le taux.
\end{definition}

Le dollar de Hong Kong est lié au dollar américain depuis le 17 octobre 1983~; son
autorité monétaire le maintient dans une bande de 7.75 à 7.85 par dollar, grâce à
une caisse d'émission, qui est prête à échanger les deux monnaies aux bornes de la
bande. Une bande comme celle du renminbi est gérée par le taux central quotidien
et par des interventions lorsque le taux approche de ses limites. Un
plancher unilatéral comme celui du franc suisse oblige la banque centrale à vendre sa propre
monnaie sans limite, ce qu'elle peut toujours faire, puisqu'elle peut la créer~; un plafond
l'oblige à vendre des devises étrangères, ce qu'elle ne peut faire que tant que ses
réserves durent.

\begin{dated}{2026-09}{Ancrages, bandes et fixings}\label{dat:m2:emerging-market-fx-and-ndfs:regimes}
Dollar de Hong Kong~: HK\$7.75 à 7.85 par dollar américain dans le cadre du Linked Exchange Rate
System, en vigueur depuis le 17 octobre 1983. Renminbi~: taux au comptant onshore dans un écart de 2~\% du
taux central, la largeur fixée depuis le 17 mars 2014~; fixing NDF CNY SAEC vers
9 h 15 à Pékin, fixing offshore CNY CNHHK vers 11 h 30 à Hong Kong
(conditions standard de juin 2023). Franc suisse~: aucun taux minimal depuis le 15 janvier
2015.
\end{dated}

Pour un trader, un ancrage crédible rend le taux presque sans risque et ses options
presque sans valeur, et c'est là le piège~: la volatilité implicite à l'ancrage mesure
la confiance du marché dans la politique, et non l'ampleur du mouvement si la politique
prend fin. Le mouvement, quand il survient, est un saut.

\section{La sortie du plancher sur le franc}

\begin{omfigure}
\begin{tikzpicture}
\begin{axis}[width=11cm,height=5cm,
  ylabel={Francs suisses pour un euro},
  tick label style={font=\small},label style={font=\small},
  xmin=2010,xmax=2017,ymin=0.95,ymax=1.5,grid=major,grid style={gray!25},
  xtick={2010,2011,2012,2013,2014,2015,2016,2017},xticklabel style={/pgf/number format/1000 sep={}},
  table/col sep=comma]
\addplot[omThm,thick] table[x=t,y=eurchf]{figdata/markets-2/18-emerging-market-fx-and-ndfs/eurchf.csv};
\draw[omProp,dashed] (axis cs:2011.68,1.2) -- (axis cs:2015.04,1.2);
\node[omProp,font=\footnotesize,anchor=north,align=center] at (axis cs:2013.3,1.185) {plancher à 1.20\\6 sept. 2011 au 15 janv. 2015};
\end{axis}
\end{tikzpicture}

\omcaption{L'euro en francs suisses, taux de référence quotidiens de la BCE, de 2010 à 2016.
Après la hausse du franc en 2010--2011, la Banque nationale suisse a maintenu le taux à 1.20 ou
au-dessus pendant plus de trois ans~; le 15 janvier 2015, elle l'a laissé filer, et le
taux de référence est passé de 1.2010 à 1.0280 en une journée. Données~: ECB Data Portal,
série EXR.D.CHF.EUR.SP00.A.}\label{fig:m2:emerging-market-fx-and-ndfs:eurchf}
\end{omfigure}

Le plancher suisse était unilatéral~: la Banque nationale suisse s'était engagée à vendre des francs
à 1.20 sans limite, et pendant trois ans elle a acheté des euros chaque fois que le marché faisait
baisser le taux. Un trader pouvait considérer le plancher comme un mur~: acheter des euros
près de 1.20, avec un ordre stop un peu en dessous, semblait être une transaction à faible
risque de perte, et pouvait être conservée avec effet de levier. Lorsque la banque centrale
s'est retirée, le 15 janvier 2015, et a abaissé dans le même temps son taux de dépôt à $-0.75\%$, le taux
n'a pas traversé les niveaux intermédiaires~; il a sauté. Un ordre stop devient un ordre au marché
lorsque son niveau est traversé par les cours, et il est exécuté au prochain cours disponible, qui
ce jour-là pouvait se trouver très loin en dessous. Les clients à effet de levier pouvaient perdre
plus que leurs dépôts, et leurs courtiers le reste.

\begin{proposition}[Effet de levier et saut]\label{prop:m2:emerging-market-fx-and-ndfs:leverage}
Une position de $L$ fois la marge déposée perd $L\,m$ fois la marge lorsque le cours saute de $m$
à son encontre avant que tout stop ne puisse être exécuté~; le solde du client
devient négatif lorsque $m > 1/L$, et l'excédent est dû au courtier.
\end{proposition}

\begin{proof}
Avec une marge $M$ et une position $LM$, un mouvement relatif $m$ fait perdre $LMm$, ce qui dépasse
$M$ exactement lorsque $Lm > 1$.
\end{proof}

Un saut de 14.4~\% transforme un effet de levier de 20 pour 1 en une perte de 2.88 fois la marge, et
un effet de levier de 50 pour 1 en une perte de 7.2 fois (voir la \cref{fig:m2:emerging-market-fx-and-ndfs:leverage}). Le
courtier de l'affaire, FXCM, a annoncé le 19 janvier 2015 que les soldes débiteurs de ses clients
avaient contraint ses filiales réglementées à lever des capitaux, et que la société
avait emprunté USD~300 millions auprès de Leucadia le 16 janvier, pour deux ans, à un taux
initial de 10~\% par an, pouvant monter jusqu'à 17~\%, afin de les couvrir.

\begin{omfigure}
\begin{tikzpicture}
\begin{axis}[width=11cm,height=4.4cm,ybar,bar width=22pt,
  xlabel={levier (position / marge)},ylabel={perte / marge},
  tick label style={font=\small},label style={font=\small},
  xmin=-0.6,xmax=3.6,ymin=0,ymax=9,grid=major,grid style={gray!25},
  xtick={0,1,2,3},xticklabels={5,10,20,50},
  nodes near coords,every node near coord/.append style={font=\scriptsize},
  table/col sep=comma]
\addplot[fill=omProp!60,draw=omProp] table[x=k,y=multiple]{figdata/markets-2/18-emerging-market-fx-and-ndfs/leverage.csv};
\draw[omThm,thick,dashed] (axis cs:-0.6,1) -- (axis cs:3.6,1);
\node[omThm,font=\scriptsize,anchor=south] at (axis cs:0.5,1.05) {dépôt perdu};
\end{axis}
\end{tikzpicture}

\omcaption{Perte sur une position longue en EURCHF, en multiple du dépôt de marge versé,
pour le mouvement du taux de référence de la BCE, de 1.2010 le 14 janvier à 1.0280 le 15
janvier 2015, si aucun stop n'a été exécuté avant le cours de l'après-midi. Au-dessus de la
ligne pointillée, le client doit de l'argent au courtier. Données~: taux de référence de la BCE~; tutoriel
du chapitre.}\label{fig:m2:emerging-market-fx-and-ndfs:leverage}
\end{omfigure}

\section{Tutoriel~: règlement du NDF et base onshore--offshore}\label{tut:m2:emerging-market-fx-and-ndfs:ndf}

\begin{tutorial}
\textbf{Objectif.} Régler un NDF par rapport à son fixing, dater le fixing, calculer les
taux d'intérêt implicites des termes onshore et offshore et leur base, ainsi que la
perte d'une position à effet de levier lors d'un saut. \textbf{État final~:}
\cref{ex:m2:emerging-market-fx-and-ndfs:basis,ex:m2:emerging-market-fx-and-ndfs:krw},
\cref{fig:m2:emerging-market-fx-and-ndfs:eurchf,fig:m2:emerging-market-fx-and-ndfs:leverage}
et les chiffres du problème du week-end.
\begin{tutsteps}
\item \textbf{Le NDF}~: son règlement et sa date de fixing.
\omcode{code/firm/ndf/firm_ndf.py}{11}{24}{Règlement du NDF et date de fixing.}\label{lst:m2:emerging-market-fx-and-ndfs:ndf}
\item \textbf{La base}~: taux locaux implicites à partir du terme de chaque marché.
\omcode{code/firm/ndf/firm_ndf.py}{27}{36}{Taux local implicite et base onshore--offshore.}\label{lst:m2:emerging-market-fx-and-ndfs:basis}
\item \textbf{Exécuter} \texttt{ndf\_demo.krw\_ndf()}, \texttt{ndf\_demo.cny\_cnh()},
\texttt{ndf\_demo.unpeg()} et \texttt{fig\_ndf.py}.
\end{tutsteps}
\textbf{Pour aller plus loin.} Évaluer la valeur du NDF avant fixing, à partir de
la courbe à terme des NDF et de l'actualisation en dollars~; puis calculer comment la base de
l'\cref{ex:m2:emerging-market-fx-and-ndfs:basis} modifie la valeur d'une couverture à trois mois
pour un détenteur étranger d'obligations chinoises.
\end{tutorial}

\section{Construction~: le calculateur de NDF}\label{bld:m2:emerging-market-fx-and-ndfs:ndf}

\begin{build}
\bfield{Rôle} L'entreprise miniature négocie des devises des marchés émergents offshore~:
elle a besoin des montants de règlement des NDF et des dates de fixing, des taux offshore implicites
des courbes de NDF, et d'une mesure du coût du levier lors d'un saut. \bfield{Interface} \texttt{ndf\_settlement(notional\_usd, contract, fixing,
buyer)}~; \texttt{fixing\_date(settlement, holidays, lag)}~;
\texttt{implied\_local\_rate}~; \texttt{onshore\_offshore\_basis}~;
\texttt{band(central, width)}~; \texttt{leveraged\_loss(move, leverage)}.
\bfield{Règles} Les taux sont exprimés en devise locale par dollar~; règlement en dollars~;
fixing deux jours ouvrés avant le règlement par défaut~; taux en dollars sur base
exact/360 et taux locaux sur base exact/365, sauf indication contraire.
\bfield{Tests d'acceptation} \texttt{code/firm/ndf/tests/}~: règlement nul au
taux du contrat et signes de part et d'autre~; une date de fixing qui saute un week-end et un
jour férié~; un taux implicite retrouvé à partir d'un terme construit avec lui~; les limites de la
bande~; le multiple de levier.
\bfield{Pour aller plus loin} Solutions de repli lorsque la \omterm{def:m2:emerging-market-fx-and-ndfs:ndf}{source du fixing} n'est pas publiée (les
règles de perturbation des conditions standard)~; courbes de NDF bootstrappées à partir de cotations~;
la valeur d'une option sur la rupture d'un \omterm{def:m2:emerging-market-fx-and-ndfs:peg}{ancrage de change}, évaluée avec un saut.
\end{build}

\begin{omsources}
\item Swiss National Bank, communiqués de presse du 6 septembre 2011 et du 15 janvier 2015~;
European Central Bank, taux de référence en euros (Data Portal).
\item FXCM Inc., Form 8-K, 19 janvier 2015 (Exhibit 99.1).
\item People's Bank of China, annonce sur la bande de négociation du renminbi, mars
2014~; Hong Kong Monetary Authority, Linked Exchange Rate System, et renminbi
business booklet.
\item EMTA et ISDA, annexe A aux 1998 FX and Currency Option Definitions, juin
2023.
\end{omsources}

\section{Exercices}

\begin{exercise}[$\star$]\label{exo:m2:emerging-market-fx-and-ndfs:1}
Un vendeur de dollars sur le NDF du won de l'\cref{ex:m2:emerging-market-fx-and-ndfs:krw}~:
que reçoit-il ou que paie-t-il à un fixing de 1\,425.50, et à 1\,350.00~?
\end{exercise}

\begin{exercise}[$\star$]\label{exo:m2:emerging-market-fx-and-ndfs:2}
Pourquoi un fonds étranger couvre-t-il des obligations en roupies indiennes par un NDF plutôt
que par un terme livrable~?
\end{exercise}

\begin{exercise}[$\star$]\label{exo:m2:emerging-market-fx-and-ndfs:3}
La parité centrale du renminbi est de 7.1000. Donner la bande de négociation onshore, ainsi que la
bande du dollar de Hong Kong.
\end{exercise}

\begin{exercise}[$\star\star$]\label{exo:m2:emerging-market-fx-and-ndfs:4}
Un NDF se règle le lundi 5 octobre 2026 et le vendredi 2 octobre est un jour férié pour sa
\omterm{def:m2:emerging-market-fx-and-ndfs:ndf}{source du fixing}. À quelle date le fixing est-il pris~?
\end{exercise}

\begin{exercise}[$\star\star$]\label{exo:m2:emerging-market-fx-and-ndfs:5}
Dans l'\cref{ex:m2:emerging-market-fx-and-ndfs:basis}, que signifie une base positive
offshore moins onshore pour un investisseur étranger qui couvre offshore des
obligations en renminbi~?
\end{exercise}

\begin{exercise}[$\star\star$]\label{exo:m2:emerging-market-fx-and-ndfs:6}
Pourquoi un plancher, comme celui du franc suisse, est-il plus facile à défendre qu'un plafond~?
\end{exercise}

\begin{exercise}[$\star\star\star$]\label{exo:m2:emerging-market-fx-and-ndfs:7}
\emph{Programmation.} Avec les données de la BCE, donner le plus bas cours de référence EURCHF
entre le 7 septembre 2011 et le 14 janvier 2015, les cours de référence des 14 et 15 janvier
2015, et le plus bas cours de janvier 2015.
\end{exercise}

\begin{exercise}[$\star\star\star$]\label{exo:m2:emerging-market-fx-and-ndfs:8}
\emph{Trouver l'erreur.} «~La volatilité implicite à un mois sur l'EURCHF vaut 1~\% avec
le plancher en place~: une position longue sur EURCHF assortie d'un stop 2~\% plus bas risque au plus 2~\%.~»
Corriger-la.
\end{exercise}

\section{Problème~: le plancher qui a cédé}

\begin{problem}[{Problème du week-end --- un client à effet de levier le 15 janvier 2015}]\label{pb:m2:emerging-market-fx-and-ndfs:1}
Un client de détail détient une position longue de 1~million d'euros contre francs à 1.2010, cours
de référence du 14 janvier 2015, avec un effet de levier de 20 pour 1 et un stop à 1.19. Le 15 janvier,
le plancher est retiré~; le cours de référence de l'après-midi est de 1.0280.

\medskip\noindent\textbf{Partie I --- La position.}
\begin{enumerate}
\item Donner la marge déposée, en francs.
\item Quelle perte le stop à 1.19 aurait-il produite~?
\item Donner la variation jusqu'à 1.0280, en pourcentage.
\item Donner la perte à 1.0280, en francs et en multiple de la marge.
\item Donner le solde du client.
\end{enumerate}

\medskip\noindent\textbf{Partie II --- Pourquoi le stop a échoué.}
\begin{enumerate}[resume]
\item Qu'est-ce qu'un ordre stop, et que garantit-il~?
\item Pourquoi n'y avait-il aucun acheteur d'euros autour de 1.19~?
\item Pourquoi le calme du plancher a-t-il fait paraître l'effet de levier sans danger~?
\item Que disait la volatilité implicite à un mois avant, et que n'a-t-elle pas perçu~?
\item Comment une option de vente sur l'EURCHF se serait-elle comportée~?
\end{enumerate}

\medskip\noindent\textbf{Partie III --- Le courtier.}
\begin{enumerate}[resume]
\item Qui supporte d'abord le solde négatif du client~?
\item Comment FXCM a-t-elle couvert les soldes négatifs de ses clients~?
\item Quel a été le coût de l'emprunt~?
\item Comment un courtier doit-il fixer l'effet de levier sur une paire soumise à un \omterm{def:m2:emerging-market-fx-and-ndfs:peg}{ancrage de change}~?
\item Qu'est-ce qui protège les clients de détail contre les soldes négatifs dans certaines juridictions~?
\end{enumerate}

\medskip\noindent\textbf{Partie IV --- Jugement.}
\begin{enumerate}[resume]
\item Pourquoi les banques centrales abandonnent-elles brutalement leurs ancrages plutôt que progressivement~?
\item Quels autres ancrages ou \omterm{def:m2:emerging-market-fx-and-ndfs:peg}{bandes de fluctuation} pourraient céder de la même façon~?
\item Comment dimensionner une position dans une devise sous \omterm{def:m2:emerging-market-fx-and-ndfs:peg}{ancrage de change}~?
\item Énoncer le \emph{résultat nommé}~: la perte de la position longue à 20 pour 1 au cours
de référence de la BCE du 15 janvier 2015.
\item En une phrase~: quel est le risque d'une devise sous \omterm{def:m2:emerging-market-fx-and-ndfs:peg}{ancrage de change}~?
\end{enumerate}
\end{problem}

\section{Questions d'entretien}

\begin{interviewq}[$\star$ \iqroles{trader, développeur}]\label{iq:m2:emerging-market-fx-and-ndfs:1}
Qu'est-ce qu'un NDF, et comment se fait son règlement~?
\end{interviewq}

\begin{interviewq}[$\star$ \iqroles{trader, chercheur}]\label{iq:m2:emerging-market-fx-and-ndfs:2}
Pourquoi le renminbi sur le \omterm{def:m2:emerging-market-fx-and-ndfs:onshore}{marché onshore} et celui sur le \omterm{def:m2:emerging-market-fx-and-ndfs:onshore}{marché offshore} peuvent-ils se négocier à des prix différents~?
\end{interviewq}

\begin{interviewq}[$\star\star$ \iqroles{chercheur}]\label{iq:m2:emerging-market-fx-and-ndfs:3}
Comment évalueriez-vous une option sur une devise dont l'ancrage pourrait céder~?
\end{interviewq}

\begin{interviewq}[$\star\star$ \iqroles{banque, trader}]\label{iq:m2:emerging-market-fx-and-ndfs:4}
Qu'est-ce qui a mal tourné pour les traders et les courtiers à effet de levier le 15 janvier 2015~?
\end{interviewq}

\begin{interviewq}[$\star\star$ \iqroles{trader}]\label{iq:m2:emerging-market-fx-and-ndfs:5}
Un client demande à couvrir 100~millions de dollars d'obligations brésiliennes. Comment structureriez-vous
l'opération, et que peut-il mal tourner au fixing~?
\end{interviewq}

\begin{interviewq}[$\star\star\star$ \iqroles{développeur, banque}]\label{iq:m2:emerging-market-fx-and-ndfs:6}
Concevoir le système de gestion des risques d'un courtier FX de détail de sorte qu'un mouvement comparable à celui du 15
janvier 2015 ne menace pas sa solvabilité.
\end{interviewq}
